\section{Related work}
\label{sec:related}

\paragraph{Cash-flow and personal-finance forecasting.} Most published
cash-flow forecasting concerns businesses and point forecasts; for example,
\citet{weytjens2021cashflow} compare neural networks with ARIMA and Prophet
on company cash flows. Work on individuals' balances is scarce, largely
because bank microdata is not public. \citet{zhang2018lowwage} built
forecasts for a budgeting app's low-wage users: a separate, small cohort
whose accounts were collected through Plaid. They also used the PKDD'99
data (Berka), which they note is not specific to low-wage workers. On
Berka they kept accounts with several years of data and aggregated
to weekly balances, because daily activity is sparse. They treated the
true balances as unknown and started every account at zero. Their method
matches the recent balance sequence against historical sequences pooled
across users, and they score point forecasts by mean absolute error. Our
results cannot be compared with theirs directly. We forecast daily
balances with predictive distributions, from one account's own history,
scored by CRPS and coverage, on a different selection of accounts.
Berka's balance field, used to validate our running balances, is
described in \S\ref{sec:protocol}. Deep probabilistic models such as DeepAR
\citep{salinas2020deepar} learn across many series and so need pooled data,
which our local-first setting does not allow; our pooled baseline
(B-GLOBAL, gradient-boosted quantile regression with LightGBM,
\citealp{ke2017lightgbm}) stands in for that family.

\paragraph{State-space models and adaptive Kalman filtering.} Structural
time-series models and their Kalman-filter treatment are standard
\citep{harvey1989,durbin2012}. Estimating unknown noise covariances
online goes back to \citet{mehra1970,mehra1972} and
\citet{myers1976}; variational Bayesian alternatives exist
\citep{sarkka2009vb}. Our filter uses covariance matching with
exponential forgetting and an information weight per sample. Page's CUSUM
\citep{page1954} is the classical sequential test for a shift in mean; we
apply it to the filter's normalized innovations.

\paragraph{Evaluating probabilistic forecasts.} We follow the paradigm of
maximizing sharpness subject to calibration
\citep{gneiting2007calibration}. Scores are the continuous ranked
probability score \citep{matheson1976,gneiting2007strictly}, the coverage
of central intervals, and the probability integral transform
\citep{dawid1984}. Sample quantiles use the type-7 rule of
\citet{hyndman1996}. Because forecasts from one account are correlated, we
compare methods by bootstrap intervals that resample whole accounts,
rather than by a test that treats forecasts as independent
\citep[cf.][]{diebold1995}.

\paragraph{Conformal prediction.} Conformal methods give distribution-free
coverage guarantees under exchangeability \citep{vovk2005alrw}. Adaptive
conformal inference \citep{gibbs2021aci} adjusts the working level online
and keeps long-run coverage under distribution shift; our variant V1
applies it per account and horizon to the Monte Carlo cone.

\paragraph{Pre-registration.} Pre-registration separates hypothesis
generation from testing \citep{nosek2018prereg}, and a NeurIPS workshop
proposed it for machine learning \citep{pmlr148prereg}. We follow a lighter
practice: an analysis plan with falsifiable claims and fixed thresholds,
kept in a version-controlled repository; a development/held-out split; a
rule for revision written before the revision; a freeze; and a held-out
evaluation that is mechanically limited to one logged run. Unlike a
registration with a third party, the repository's commit times are under
the author's control until they are pushed, which we account for in
\S\ref{sec:timeline}.
